\documentclass[11pt,a4paper]{article}

\usepackage[utf8]{inputenc}
\usepackage[T1]{fontenc}
\usepackage{amsmath,amssymb,amsthm}
\usepackage{geometry}
\usepackage{booktabs}
\usepackage{xcolor}
\usepackage{physics}
\usepackage{hyperref}
\usepackage{enumitem}
\usepackage{tikz}

\geometry{margin=1in}

\newtheorem{theorem}{Theorem}[section]
\newtheorem{proposition}[theorem]{Proposition}
\newtheorem{lemma}[theorem]{Lemma}
\newtheorem{corollary}[theorem]{Corollary}
\theoremstyle{definition}
\newtheorem{definition}[theorem]{Definition}
\newtheorem{remark}[theorem]{Remark}
\newtheorem{example}[theorem]{Example}

\newcommand{\Tc}{\mathcal{T}}
\newcommand{\dgeo}{\delta}

\title{\textbf{QSD Supplement S1: Formal Proofs}\\[0.5em]
\large Graph Spectra, Group Theory, and Algebraic Foundations}

\author{Marcus Hubka\\
\small Independent Researcher, Czech Republic}

\date{February 2026}

\begin{document}

\maketitle

\begin{abstract}
This supplement provides rigorous mathematical proofs for the core results of Quantum Spectrodynamics. We derive the complete spectrum of the Heisenberg Hamiltonian on $K_4$, prove the emergence of gauge group structure from $S_4$ representation theory, establish the golden ratio identity from the $C_5$ Laplacian, and demonstrate the origin of the geometric strain parameter.
\end{abstract}

\tableofcontents

\section{The Heisenberg Hamiltonian on $K_4$}

\subsection{Setup and Definitions}

\begin{definition}[Heisenberg Hamiltonian on $K_4$]
\begin{equation}
H_{K_4} = J \sum_{i<j}^{4} \mathbf{S}_i \cdot \mathbf{S}_j
\end{equation}
where $J > 0$ is the exchange coupling and $\mathbf{S}_i = \frac{1}{2}\boldsymbol{\sigma}_i$ are spin-1/2 operators.
\end{definition}

The Hilbert space has dimension $2^4 = 16$.

\subsection{Total Spin Method}

\begin{lemma}[Casimir Trick]
\begin{equation}
\sum_{i<j} \mathbf{S}_i \cdot \mathbf{S}_j = \frac{1}{2}\left( \mathbf{S}_{\text{tot}}^2 - \sum_i \mathbf{S}_i^2 \right)
\end{equation}
\end{lemma}

\begin{proof}
Expand $\mathbf{S}_{\text{tot}}^2 = (\sum_i \mathbf{S}_i)^2 = \sum_i \mathbf{S}_i^2 + 2\sum_{i<j} \mathbf{S}_i \cdot \mathbf{S}_j$.
\end{proof}

For spin-1/2, $\mathbf{S}_i^2 = 3/4$. Thus:
\begin{equation}
H_{K_4} = \frac{J}{2}\left( \mathbf{S}_{\text{tot}}^2 - 3 \right)
\end{equation}

\subsection{Spin Coupling}

\begin{proposition}[Spin Coupling]
Coupling 4 spin-1/2 particles yields $S_{\text{tot}} \in \{0, 1, 2\}$ with multiplicities 2, 3, 1 respectively.
\end{proposition}

\begin{proof}
Using $\frac{1}{2} \otimes \frac{1}{2} = 0 \oplus 1$ twice, then combining intermediate spins. Dimension check: $2(1) + 3(3) + 1(5) = 16$ \checkmark
\end{proof}

\begin{theorem}[$K_4$ Spectrum]
\begin{align}
E(S=0) &= -\frac{3J}{2} \quad \text{(deg. 2)} \\
E(S=1) &= -\frac{J}{2} \quad \text{(deg. 9)} \\
E(S=2) &= +\frac{3J}{2} \quad \text{(deg. 5)}
\end{align}
\end{theorem}

\section{$S_4$ Representation Theory}

\subsection{Conjugacy Classes}

$S_4$ has 5 conjugacy classes with sizes 1, 6, 8, 6, 3.

\subsection{Irreducible Representations}

$S_4$ has irreps $\mathbf{1}, \mathbf{1}', \mathbf{2}, \mathbf{3}, \mathbf{3}'$ with dimensions 1, 1, 2, 3, 3.

\begin{theorem}[Edge Decomposition]
The 6 edges of $K_4$ transform as:
\begin{equation}
\mathbf{6}_{\text{edges}} = \mathbf{1} \oplus \mathbf{2} \oplus \mathbf{3}
\end{equation}
\end{theorem}

\begin{proof}
Compute characters: $\chi = (6, 2, 0, 0, 2)$. Verify $\chi_{\mathbf{1}} + \chi_{\mathbf{2}} + \chi_{\mathbf{3}} = \chi$.
\end{proof}

This maps to $U(1)_Y \times SU(2)_L \times SU(3)_c$.

\section{The $C_5$ Laplacian}

\subsection{Eigenvalues}

\begin{theorem}[$C_5$ Laplacian Spectrum]
\begin{equation}
\lambda_k = 2 - 2\cos\left(\frac{2\pi k}{5}\right), \quad k = 0,1,2,3,4
\end{equation}
giving $\lambda = \{0, 1.382, 3.618, 3.618, 1.382\}$.
\end{theorem}

\subsection{Golden Ratio}

\begin{theorem}[Golden Ratio Emergence]
\begin{equation}
\frac{\lambda_2}{\lambda_1} = \frac{5 + \sqrt{5}}{5 - \sqrt{5}} = \varphi^2 \approx 2.618
\end{equation}
\end{theorem}

\begin{proof}
Use exact values $\cos(72°) = (\sqrt{5}-1)/4$ and $\cos(144°) = -(\sqrt{5}+1)/4$.
\end{proof}

\section{The Geometric Strain}

\begin{definition}[Geometric Strain]
\begin{equation}
\dgeo = \frac{3}{80\varphi^2} = 0.01432
\end{equation}
\end{definition}

This measures the incompatibility between $K_4$ (6 edges) and $C_5$ (5 edges).

\end{document}
